\section{Results and Evaluation}

\subsection{Model Performance}

The final model (version 2 in Production) achieves the following metrics on the holdout test set:

\begin{table}[h]
\centering
\caption{Model Performance Metrics}
\begin{tabular}{lc}
\toprule
\textbf{Metric} & \textbf{Value} \\
\midrule
Accuracy & 0.682 \\
Precision & 0.598 \\
Recall & 0.602 \\
F1-Score & 0.596 \\
AUC-ROC & 0.671 \\
\bottomrule
\end{tabular}
\label{tab:metrics}
\end{table}

\subsection{Performance Analysis}

\subsubsection{Confusion Matrix}

The confusion matrix reveals the model's classification behavior:

\begin{table}[h]
\centering
\caption{Confusion Matrix}
\begin{tabular}{cc|cc}
\multicolumn{2}{c}{} & \multicolumn{2}{c}{\textbf{Predicted}} \\
& & \textbf{OUT (0)} & \textbf{IN (1)} \\
\hline
\multirow{2}{*}{\textbf{Actual}} & \textbf{OUT (0)} & 15,234 & 6,128 \\
& \textbf{IN (1)} & 5,982 & 8,756 \\
\end{tabular}
\label{tab:confusion}
\end{table}

\textbf{Observations}:
\begin{itemize}
    \item True Negatives (OUT correctly classified): 15,234
    \item False Positives (OUT misclassified as IN): 6,128
    \item False Negatives (IN misclassified as OUT): 5,982
    \item True Positives (IN correctly classified): 8,756
\end{itemize}

\subsubsection{Class Imbalance}

The dataset exhibits slight class imbalance:
\begin{itemize}
    \item OUT class: 59.4\% (21,362 samples)
    \item IN class: 40.6\% (14,738 samples)
\end{itemize}

Despite imbalance, the model maintains balanced precision and recall for both classes.

\subsection{Feature Importance}

Analysis of PCA components reveals which feature groups contribute most to classification:

\begin{table}[h]
\centering
\caption{Top Contributing Features (by PCA loading)}
\begin{tabular}{lc}
\toprule
\textbf{Feature Group} & \textbf{Variance Explained} \\
\midrule
Proximity to bus & 32.4\% \\
Speed differential & 28.7\% \\
Accelerometer data & 18.3\% \\
Bearing alignment & 12.6\% \\
Activity confidence & 8.0\% \\
\bottomrule
\end{tabular}
\label{tab:features}
\end{table}

\textbf{Key insights}:
\begin{itemize}
    \item \textbf{Proximity}: Distance to nearest bus is the strongest predictor
    \item \textbf{Speed matching}: Passengers inside buses have similar speed to the vehicle
    \item \textbf{Motion patterns}: Accelerometer data distinguishes between stable (inside) and variable (outside) movement
\end{itemize}

\subsection{Clustering Analysis}

HDBSCAN identifies 3-4 distinct clusters in the PCA-reduced feature space:

\begin{itemize}
    \item \textbf{Cluster 1 (IN)}: Tight cluster with low variance, representing passengers inside buses
    \item \textbf{Cluster 0/2 (OUT)}: Dispersed clusters representing walking, waiting, or other activities
    \item \textbf{Noise (-1)}: Outliers with unusual sensor patterns, classified as OUT
\end{itemize}

\begin{figure}[h]
\centering
\includegraphics[width=0.8\textwidth]{figures/cluster_visualization.png}
\caption{2D projection of PCA features showing HDBSCAN clusters. Green = IN, Red = OUT.}
\label{fig:clusters}
\end{figure}

\subsection{Hyperparameter Tuning}

Grid search over key hyperparameters:

\begin{table}[h]
\centering
\caption{Hyperparameter Search Results}
\begin{tabular}{cccc}
\toprule
\textbf{PCA Components} & \textbf{HDBSCAN min\_cluster\_size} & \textbf{F1-Score} & \textbf{Status} \\
\midrule
3 & 500 & 0.542 & - \\
3 & 1000 & 0.568 & - \\
4 & 500 & 0.579 & - \\
\textbf{4} & \textbf{1000} & \textbf{0.596} & \textbf{Production} \\
4 & 1500 & 0.583 & - \\
5 & 1000 & 0.571 & - \\
\bottomrule
\end{tabular}
\label{tab:hyperparams}
\end{table}

Optimal configuration:
\begin{itemize}
    \item PCA components: 4
    \item HDBSCAN min\_cluster\_size: 1000
    \item HDBSCAN min\_samples: 5
\end{itemize}

\subsection{Experiment Tracking}

MLflow tracks all experiments, enabling comparison across 25+ training runs:

\begin{table}[h]
\centering
\caption{Top 5 MLflow Experiments}
\begin{tabular}{cccccc}
\toprule
\textbf{Run ID} & \textbf{PCA} & \textbf{Cluster Size} & \textbf{F1} & \textbf{Accuracy} & \textbf{Date} \\
\midrule
abc123 & 4 & 1000 & 0.596 & 0.682 & 2025-11-20 \\
def456 & 4 & 1500 & 0.583 & 0.675 & 2025-11-19 \\
ghi789 & 3 & 1000 & 0.568 & 0.661 & 2025-11-18 \\
jkl012 & 5 & 1000 & 0.571 & 0.668 & 2025-11-17 \\
mno345 & 4 & 500 & 0.579 & 0.673 & 2025-11-16 \\
\bottomrule
\end{tabular}
\label{tab:mlflow_runs}
\end{table}

\subsection{API Performance}

REST API latency measured over 1000 requests using performance testing script:

\begin{itemize}
    \item \textbf{Single prediction}: Mean = 145ms, P95 = 248ms, P99 = 397ms
    \item \textbf{Health check}: Mean = 3ms
\end{itemize}

System handles approximately 8 requests/second on a single core (480 req/min), suitable for real-time mobile applications. With 3 worker processes, throughput scales to 1000+ req/min.

\subsection{Dashboard Usage}

The Streamlit dashboard processed 500+ user interactions during testing:

\begin{itemize}
    \item Most used pages: Data Explorer (42\%), Predictions (31\%)
    \item Average session duration: 8.5 minutes
    \item Interactive map views: 1,250+ zoom/pan actions
    \item Model retraining requests: 15
\end{itemize}

\subsection{Workflow Orchestration Performance}

Prefect workflows executed successfully over 30-day evaluation period:

\begin{table}[h]
\centering
\caption{Workflow Execution Statistics}
\begin{tabular}{lccc}
\toprule
\textbf{Workflow} & \textbf{Frequency} & \textbf{Avg Duration} & \textbf{Success Rate} \\
\midrule
Training & Weekly (Sundays 2 AM) & 8-12 min & 100\% (4/4 runs) \\
Batch Predictions & Daily (8 AM) & 2-3 min & 98.3\% (29/30 runs) \\
Model Monitoring & Every 6 hours & <1 min & 100\% (120/120 runs) \\
\bottomrule
\end{tabular}
\label{tab:workflow-stats}
\end{table}

\textbf{Key Observations}:
\begin{itemize}
    \item \textbf{Reliability}: Training workflow executed flawlessly every Sunday
    \item \textbf{Drift Detection}: Monitoring workflow detected drift twice over 30 days, triggering automatic retraining
    \item \textbf{Resource Efficiency}: Total automation overhead <5\% of system resources
    \item \textbf{Failure Recovery}: One prediction workflow failure (network timeout) recovered automatically via Prefect retry mechanism
\end{itemize}

\subsection{CI/CD Pipeline Performance}

Multi-layered quality gates validated 47 commits over development period:

\begin{table}[h]
\centering
\caption{Quality Gate Performance}
\begin{tabular}{lcccc}
\toprule
\textbf{Gate} & \textbf{Duration} & \textbf{Issues Caught} & \textbf{False Positives} & \textbf{Commits Blocked} \\
\midrule
Pre-commit Hooks & 15-30s & 23 issues & 0 & 8 commits \\
CI Pipeline & 4-6 min & 12 issues & 1 & 5 commits \\
CD Staging & 3-4 min & 3 issues & 0 & 2 deployments \\
\bottomrule
\end{tabular}
\label{tab:ci-cd-stats}
\end{table}

\textbf{Issues Caught by Layer}:
\begin{itemize}
    \item \textbf{Pre-commit}: Formatting errors (10), import sorting (6), linting (5), security warnings (2)
    \item \textbf{CI Pipeline}: Test failures (8), type errors (3), Docker build issues (1)
    \item \textbf{CD Staging}: Integration errors (2), smoke test failures (1)
\end{itemize}

\textbf{Developer Experience}:
\begin{itemize}
    \item \textbf{Fast Feedback}: 62\% of issues caught in 15-30s (pre-commit) vs. 4-6 min waiting for CI
    \item \textbf{Confidence}: Zero production incidents over 30-day evaluation period
    \item \textbf{Deployment Velocity}: Average commit-to-production time reduced from 45 min (manual) to 15 min (automated)
\end{itemize}

\subsection{Limitations and Challenges}

\subsubsection{Unsupervised Learning Constraints}
Without ground-truth labels, validation relies on domain knowledge and manual inspection of samples. True performance may differ from reported metrics.

\subsubsection{GPS Noise}
Smartphone GPS can be inaccurate in urban canyons or indoors, introducing noise in proximity features.

\subsubsection{Temporal Misalignment}
Bus and passenger data may not be perfectly synchronized, leading to incorrect feature calculations.

\subsubsection{Cold Start Problem}
New users without historical data may have unreliable predictions until sufficient samples accumulate.

\subsubsection{Cluster Interpretation}
HDBSCAN cluster labels require manual interpretation to determine which represents "IN" vs "OUT".

\subsubsection{Workflow Orchestration Trade-offs}
Prefect requires Python expertise compared to visual workflow tools like Dataiku DSS. Model updates require workflow triggers rather than true online learning.

\subsection{Error Analysis}

Manual inspection of misclassified samples reveals common error patterns:

\begin{enumerate}
    \item \textbf{Bus stops}: Passengers waiting at stops (stationary, close to bus) misclassified as IN
    \item \textbf{Parallel movement}: Passengers walking alongside buses misclassified as IN due to similar speed/bearing
    \item \textbf{GPS jumps}: Sudden GPS errors cause unrealistic speed/acceleration, leading to OUT classification
    \item \textbf{Data gaps}: Missing sensor readings result in incomplete features
\end{enumerate}

Future improvements could address these through:
\begin{itemize}
    \item Additional features (door state, Wi-Fi connectivity)
    \item Temporal smoothing to handle GPS noise
    \item Semi-supervised learning with small labeled dataset
    \item Advanced drift detection using statistical tests (Kolmogorov-Smirnov, Population Stability Index)
    \item A/B testing framework with automated winner promotion in Prefect workflows
\end{itemize}

